% Reference Document Utility Report
% Skill Development & Career Readiness for FYE Student Employees

\documentclass[11pt,letterpaper]{article}

\usepackage[margin=1in]{geometry}
\usepackage{helvet}
\renewcommand{\familydefault}{\sfdefault}
\usepackage[T1]{fontenc}
\usepackage[utf8]{inputenc}
\usepackage{array}
\usepackage{tabularx}
\usepackage{booktabs}
\usepackage{url}

\pagestyle{plain}

\setlength{\parindent}{0pt}
\setlength{\parskip}{0.45em}

\makeatletter
\renewcommand\section{\@startsection{section}{1}{\z@}%
  {1.1em}{0.45em}{\large\bfseries}}
\renewcommand\subsection{\@startsection{subsection}{2}{\z@}%
  {0.85em}{0.3em}{\normalsize\bfseries}}
\makeatother

\newcommand{\sectionrule}{\par\vspace{0.2em}\noindent\rule{\textwidth}{0.5pt}\vspace{0.5em}}

\begin{document}

\begin{center}
  {\LARGE\bfseries Reference Document Utility Report}\\[0.35em]
  {\large What Each Source Contains and How It Supports the Project}\\[0.6em]
  {\normalsize Skill Development \& Career Readiness for FYE Student Employees}\\[0.8em]
  {\small Prepared for: FYE Research \& Engagement Project}\\
  {\small Student: Pushpak Sunil Rane \quad|\quad Supervisor: Dani Fernandez}\\
  {\small California State University, Chico \quad|\quad Fall 2026}
\end{center}

\sectionrule

\begin{abstract}
\noindent
This report inventories the research reference documents used for the FYE study and explains, for each source: (1)~what it contains, (2)~how it can be used in the FYE study, and (3)~how it connects to the project's goals. Project-management files (such as the REC proposal template and internal utility summary) are excluded here; they are used only for planning and updates, not as research evidence. The project assesses skill development and career readiness practices built into FYE student employment---weekly competency reflection, assessment and feedback, guided professional goal-setting, and interview support---across new employees, returning employees, and past employees in post-graduate careers. Together, the research references provide local practice benchmarks, assessment templates, career-education tools, and external employer/future-skills evidence needed to design, justify, and improve the study.
\end{abstract}

\sectionrule

\section{Project Connection (What This Report Must Support)}
The FYE Research \& Engagement project seeks to understand and strengthen how student employment builds career readiness. Specifically, the study focuses on:
\begin{enumerate}
  \item Weekly competency reflection;
  \item Assessment and feedback (self-report + supervisor/lead feedback);
  \item Guided professional goal-setting;
  \item Interview support;
  \item Multi-cohort impact (new, returning, and alumni/post-graduate employees);
  \item Program improvement recommendations for FYE.
\end{enumerate}

The research references fall into two groups:
\begin{itemize}
  \item \textbf{Local / campus practice documents} that describe how student employment and evaluation currently work;
  \item \textbf{External frameworks and labor-market evidence} that define which skills matter for career readiness now and through 2030.
\end{itemize}
Files used only for project updating (REC proposal template; internal markdown utility summary) are not treated as research references in this report.

\section{Document Inventory at a Glance}

\begin{center}
\renewcommand{\arraystretch}{1.15}
\begin{tabularx}{\textwidth}{>{\raggedright\arraybackslash}p{0.34\textwidth} >{\raggedright\arraybackslash}p{0.22\textwidth} >{\raggedright\arraybackslash}X}
\toprule
\textbf{Document} & \textbf{Type} & \textbf{Primary role for this project} \\
\midrule
Master Student Employment Agreement.docx & Campus employment norms & Benchmark for professionalism, check-ins, and developmental framing \\
\midrule
Student Assistant Evaluation Sheet.xlsx & Evaluation instrument & Template for rubrics, mid/end assessment, reflective prompts \\
\midrule
career-handbook-2.pdf & Career education toolkit & NACE competencies, skill statements, STAR interview tools \\
\midrule
2026-job-outlook-report-spring-update.pdf & NACE employer survey & Current hiring/skills demand; skills-with-examples finding \\
\midrule
WEF\_Future\_of\_Jobs\_Report\_2025.pdf & Global skills outlook & Future skill taxonomy (resilience, leadership, AI literacy) \\
\bottomrule
\end{tabularx}
\end{center}

\section{Document-by-Document Analysis}

\subsection{Master Student Employment Agreement (Career Center)}
\paragraph{What it has.}
This document is the Career Center Student Employment Agreement. It sets expectations for student assistants, including:
\begin{itemize}
  \item Time etiquette, absences, customer service, and professional conduct;
  \item Confidentiality and FERPA responsibilities;
  \item Dress code and modeling professionalism for other students;
  \item Timesheet accuracy and employment conditions;
  \item Mandatory monthly staff meetings/trainings;
  \item On-the-job and peer training;
  \item Mid-semester and end-of-semester check-ins;
  \item Progressive response to unmet expectations.
\end{itemize}
Importantly, it frames student employment as developmental: helping students prepare for post-college careers, not only complete office tasks.

\paragraph{How it can be used.}
\begin{itemize}
  \item Campus benchmark for comparing Career Center vs.\ FYE developmental systems;
  \item Shared language for professionalism, confidentiality, and performance management;
  \item Inspiration for mid/end check-in questions in supervisor or student interviews;
  \item Local-context evidence that Chico State already treats student jobs as career preparation.
\end{itemize}

\paragraph{Project connection.}
Supports the project's claim that campus student employment is already intended as career readiness practice. Useful when explaining why FYE assessment of reflection, feedback, and professionalism is aligned with broader campus norms.

\paragraph{Limitation.}
Career Center--specific (not FYE's exact protocol). Evaluation categories lean toward front-desk/customer-service roles.

\subsection{Student Assistant Evaluation Sheet}
\paragraph{What it has.}
An Excel supervisor evaluation form using a 1--4 scale:
\begin{quote}
1 = Unsatisfactory;\quad 2 = Needs Improvement;\quad 3 = Meets Expectations;\quad 4 = Exceeds Expectations
\end{quote}
Rated domains include:
\begin{itemize}
  \item Professionalism (communication, confidentiality, attire, staff interaction);
  \item Customer Service (greeting, demeanor, patience, phone etiquette);
  \item Performance (accuracy, initiative, following directions, workflow proficiency);
  \item Work Habits (attendance, punctuality, productivity, shift handoff).
\end{itemize}
It also includes supervisor comments, dual signatures, and student reflective prompts such as:
\begin{itemize}
  \item What do you like best about your job?
  \item What challenges you most?
  \item Are there skills you would like to use and/or develop?
\end{itemize}

\paragraph{How it can be used.}
\begin{itemize}
  \item Baseline instrument anatomy for FYE mid/end evaluation redesign;
  \item Keep the 1--4 scale while remapping items to FYE competencies;
  \item Mixed-methods design: quantitative ratings + open-ended skill-development prompts;
  \item Goal-setting link through the ``skills to develop'' question;
  \item Alumni adaptation (``skills used after graduation,'' ``skills wished developed more'').
\end{itemize}

\paragraph{Project connection.}
Directly supports assessment/feedback and guided goal-setting components. Provides a ready structure for comparing new vs.\ returning employees and for building more FYE-relevant rubrics.

\paragraph{Limitation.}
Tuned to Career Center/SEO workflow. Phone/employer items should be replaced with FYE-relevant behaviors (research support, peer coaching, facilitation, competency feedback quality).

\subsection{Career Handbook (Chico State Career Center, Vol.~25)}
\paragraph{What it has.}
A 26-page handbook covering job search, LinkedIn, AI use in career development, employer-sought skills, resumes/CVs, cover letters, career fair/elevator pitch, and interviewing. Highest-value content for this project includes:
\begin{itemize}
  \item NACE Career Readiness Competencies: Communication, Critical Thinking, Teamwork, Professionalism, Equity \& Inclusion, Technology, Career \& Self-Development, Leadership;
  \item Skill-statement method: Power/Action Verb + Transferable Skill + Task/Result (basic $\rightarrow$ advanced $\rightarrow$ advanced with results);
  \item Interviewing guidelines and STAR method (Situation--Task--Action--Result);
  \item Sample behavioral questions and before/during/after interview checklists;
  \item AI prompts for interview practice and career action planning.
\end{itemize}

\paragraph{How it can be used.}
\begin{itemize}
  \item Local teaching framework to align FYE competencies with national career readiness language;
  \item Curriculum for interview-support workshops (huddles, UNIV~389, mid-semester PD);
  \item Weekly reflection scaffold (one STAR story per targeted competency);
  \item Coaching script for Training Leads giving competency feedback;
  \item Alumni survey constructs (which skills mattered most after graduation).
\end{itemize}

\paragraph{Project connection.}
This is the strongest practical toolkit for the project's interview-support and evidence-based skill-story goals. It also supplies the NACE competency language used in the project's alignment slides and recommendations.

\paragraph{Limitation.}
General career education resource; not FYE outcome data. Use as intervention content and framework, not as proof of FYE impact.

\subsection{NACE Job Outlook 2026 Spring Update}
\paragraph{What it has.}
A National Association of Colleges and Employers (NACE) survey update (fielded February~12--March~17, 2026; roughly 185 employer respondents) on Class of 2026 hiring, AI, and skills/attributes employers seek. Key findings include:
\begin{itemize}
  \item Employers project about a 5.6\% hiring increase for the Class of 2026;
  \item High-priority resume attributes include teamwork, communication, problem-solving, initiative, analytical skills, adaptability, leadership potential, detail orientation, and technical/computer skills;
  \item Employers strongly prefer candidates who demonstrate skills \textbf{with examples}, not merely list them;
  \item AI skill demand is rising for entry-level work (tool selection, prompting, productivity use, reviewing AI outputs).
\end{itemize}

\paragraph{How it can be used.}
\begin{itemize}
  \item Rationale for why weekly reflection must produce evidence-based narratives;
  \item Benchmark FYE-emphasized skills against national employer priorities;
  \item Justify interview/STAR practice as career readiness, not optional soft training;
  \item Support an optional AI/tech literacy recommendation for FYE;
  \item Inform alumni questions about whether graduates can show FYE-based examples in hiring.
\end{itemize}

\paragraph{Project connection.}
Validates the project's core improvement area: students already do meaningful FYE work; they need practice proving those skills. This source strengthens literature/rationale sections and recommendation briefs for supervisors/admins.

\paragraph{Limitation.}
National sample, not Chico State--specific. Verify exact chart percentages when citing formally.

\subsection{WEF Future of Jobs Report 2025}
\paragraph{What it has.}
World Economic Forum insight report (January 2025) based on 1,000+ employers across industries and economies, projecting 2025--2030 skill and job change. Core skills emphasized include analytical thinking; resilience/flexibility/agility; leadership and social influence; creative thinking; motivation and self-awareness; technological literacy; empathy and active listening; curiosity and lifelong learning. Fastest-growing skills include AI and big data, cybersecurity, technological literacy, creative thinking, resilience, lifelong learning, and leadership. The report also highlights large-scale upskilling needs and skill gaps as a major transformation barrier.

\paragraph{How it can be used.}
\begin{itemize}
  \item Future-proof FYE competency language (resilience, lifelong learning, leadership, tech/AI);
  \item Justify weekly reflection as lifelong-learning practice, not only semester compliance;
  \item Frame Training Lead mentoring as early talent-management / leadership practice;
  \item Support discussion sections that place FYE in global upskilling discourse;
  \item Inform alumni constructs about preparedness for skill change.
\end{itemize}

\paragraph{Project connection.}
Provides the long-horizon skill taxonomy that complements NACE's nearer-term employer view. Strengthens the argument that FYE should build durable human skills plus emerging tech literacy.

\paragraph{Limitation.}
Global/macro source. Use for foresight and framing, not as local outcome evidence.

\section{How the Documents Work Together}

\begin{center}
\renewcommand{\arraystretch}{1.15}
\begin{tabularx}{\textwidth}{>{\raggedright\arraybackslash}p{0.3\textwidth} >{\raggedright\arraybackslash}p{0.32\textwidth} >{\raggedright\arraybackslash}X}
\toprule
\textbf{External demand} & \textbf{Local practice} & \textbf{Project use} \\
\midrule
WEF future skills & Employment Agreement check-ins & Define which skills matter \\
NACE Job Outlook & Evaluation Sheet rubrics & Design/adapt assessment language \\
Career Handbook (NACE 8 + STAR) & FYE weekly process (mapped in slides) & Build interview + reflection supports \\
\bottomrule
\end{tabularx}
\end{center}

\vspace{0.4em}
Practical synthesis:
\begin{itemize}
  \item Use \textbf{WEF + NACE + Career Handbook} to define which competencies matter and how to evidence them;
  \item Use \textbf{Employment Agreement + Evaluation Sheet} to benchmark campus practice and adapt instruments;
  \item Combine all five research references for process mapping, skill alignment, gap analysis, and recommendations.
\end{itemize}

\section{Connection to Each Project Component}

\begin{center}
\renewcommand{\arraystretch}{1.15}
\begin{tabularx}{\textwidth}{>{\raggedright\arraybackslash}p{0.23\textwidth} >{\raggedright\arraybackslash}X >{\raggedright\arraybackslash}p{0.28\textwidth}}
\toprule
\textbf{Project component} & \textbf{Best reference supports} & \textbf{Immediate use} \\
\midrule
Weekly reflection & Handbook skill statements; WEF lifelong learning; NACE examples finding & STAR/example-based Self-Report prompts \\
\midrule
Assessment \& feedback & Evaluation Sheet; Employment Agreement check-ins; NACE evidence standard & Mid/end rubric adaptation \\
\midrule
Goal-setting & Evaluation ``skills to develop'' prompt; Handbook career action planning & 1:1 goal worksheet language \\
\midrule
Interview support & Handbook STAR section; NACE examples preference & Interview practice recommendations \\
\midrule
Multi-cohort inquiry & All frameworks + Evaluation Sheet prompts & New/returning/alumni question banks \\
\midrule
Program improvement & Combined gap analysis across all sources & Recommendations brief for FYE \\
\bottomrule
\end{tabularx}
\end{center}

\section{What Is Strong vs.\ Still Missing}

\subsection{Already strong in the research references}
\begin{itemize}
  \item Clear national/global skill frameworks (NACE, WEF);
  \item Campus career-readiness teaching tools (handbook);
  \item Local evaluation and professionalism templates;
  \item Enough material to support skill alignment, gap analysis, and recommendations.
\end{itemize}

\subsection{Important gaps still needed for full empirical assessment}
\begin{itemize}
  \item Official FYE competency list/rubric used in Work Package;
  \item Blank/sample Work Package pages (Practice Plan, Self-Report, feedback);
  \item Mid-semester evaluation and rehire/exit interview protocols;
  \item UNIV~389 syllabus and reflection assignments;
  \item De-identified artifact samples and/or survey responses;
  \item Alumni contact/consent process details.
\end{itemize}
These missing pieces are what turn the current folder from a framework library into a complete evidence dataset.

\section{Recommended Uses Going Forward}
\begin{enumerate}
  \item Keep using the references to justify recommendations (especially STAR practice, evidence-based stories, light AI/tech readiness).
  \item Cite NACE/WEF in supervisor briefs and final report rationale sections.
  \item Adapt Evaluation Sheet structure when FYE instruments are finalized.
  \item Use Career Handbook language in any interview-support recommendation or pilot.
  \item Request the missing FYE artifacts listed above to move from document analysis to outcome assessment.
\end{enumerate}

\section{Conclusion}
The research references already provide a strong foundation for the FYE career-readiness project. Local documents show how student employment and evaluation are organized on campus; external documents show what employers and future labor markets value. Connected to the project, these five sources support process mapping, skill alignment, gap analysis, and recommendations---and they prepare the ground for later instrument design and multi-cohort data collection once FYE-specific Work Package materials are available.

\section{Source Files Reviewed}
\begin{itemize}
  \item \texttt{references/Master Student Employment Agreement.docx}
  \item \texttt{references/Student Assistant Evaluation Sheet.xlsx}
  \item \texttt{references/career-handbook-2.pdf}
  \item \texttt{references/2026-job-outlook-report-spring-update.pdf}
  \item \texttt{references/WEF\_Future\_of\_Jobs\_Report\_2025.pdf}
\end{itemize}

\vspace{0.4em}
\noindent\textit{Excluded from research references (project-update files only):}\\
\texttt{F26\_ REC Project Proposal [Pushpak].pdf};\\
\texttt{Document\_Utility\_Summary\_FYE\_Career\_Readiness\_Study.md}.

\vspace{1em}
{\footnotesize\textit{Note: This report synthesizes currently available reference materials for project planning and supervisor communication. Formal citations in the final academic report should quote exact figures from the original PDFs.}}

\end{document}
